% GENERATED FILE. Edit canonical lessons, then run npm run generate:books.
% canonical-source-hash: fnv1a64:6c8bba5090e273bb
% canonical-lessons: RU-C450-spravka, RU-C450-blank, RU-C450-ochered, RU-C450-oformlyat, RU-C450-zaveryat, RU-R450-review-443-450

\chapter{Certificates and forms}
\label{ch:ru-certificates-and-forms}
\hlchaptermodality{450}

\begin{chapteropening}
\textbf{By the end of this chapter:} \emph{I can deal with paperwork: a certificate, a form, a queue, getting documents done and certifying a copy.}

Complete the last lesson of chapter 450: I can deal with paperwork: a certificate, a form, a queue, getting documents done and certifying a copy.
\end{chapteropening}

\section[\texorpdfstring{\ru{справка} (správka)}{справка (správka)}]{\ru{справка} (správka) — a certificate, an official note; information}

\label{lesson:RU-C450-spravka}



\begin{quote}
Before the new one: say the Russian for to shave, then the Russian for a fringe.
\end{quote}

\subsection*{You'll want to know: \ru{справка}}
\textbf{\ru{справка}} (\emph{správka}) — \textquotedblleft{}a certificate, an official note; information\textquotedblright{}.

\begin{quote}
\ru{Для} \ru{визы} \ru{нужна} \ru{справка} \ru{с} \ru{работы}.
\end{quote}

\begin{quote}
\emph{Dlya vízy nuzhná správka s rabóty.}
\end{quote}

\begin{quote}
For the visa you need a letter from your employer.
\end{quote}

Feminine. \emph{\ru{Справочное} \ru{бюро}} — an information desk.

\begin{cousinweb}
\emph{\ru{справка}} is built on \emph{\ru{справиться}} 'to inquire', from \emph{\ru{править}} 'to set right, to put in order'. EtymDB lists \emph{\ru{справка}} without a deeper source.
\end{cousinweb}

\begin{grammarlens}[title={What you've built}]
One more word for telling, explaining and describing at B1.
\end{grammarlens}

\subsection*{Your turn}
\begin{itemize}
\raggedright
  \item \emph{Say it:} \emph{správka}
  \item \emph{Say it:} \emph{správka}, once more
  \item \emph{Say it:} a sentence of your own with \emph{správka}
  \item \emph{Recall:} say the Russian for to shave, then the Russian for a fringe, then say \emph{správka} again
\end{itemize}

\begin{tcolorbox}[breakable,colback=teal!4,colframe=teal!35!black,title={Before you move on}]
What does \ru{справка} mean? (\textquotedblleft{}A certificate, an official note; information\textquotedblright{}.) Say it once more.
\end{tcolorbox}
\section[\texorpdfstring{\ru{бланк} (blank)}{бланк (blank)}]{\ru{бланк} (blank) — a form (to fill in)}

\label{lesson:RU-C450-blank}



\begin{quote}
Before the new one: say the Russian for a fringe, then the Russian for a certificate.
\end{quote}

\subsection*{You'll want to know: \ru{бланк}}
\textbf{\ru{бланк}} (\emph{blank}) — \textquotedblleft{}a form (to fill in)\textquotedblright{}.

\begin{quote}
\ru{В} \ru{банке} \ru{мне} \ru{дали} \ru{бланк}.
\end{quote}

\begin{quote}
\emph{V bánke mne dáli blank.}
\end{quote}

\begin{quote}
At the bank they gave me a form.
\end{quote}

Masculine. \emph{\ru{Заполнить} \ru{бланк}} — to fill in a form.

\begin{cousinweb}
\emph{\ru{бланк}} is usually traced, probably through German, to French \emph{blanc} 'white, blank' (EtymDB lists the French word): a sheet left blank to be filled in.
\end{cousinweb}

\begin{grammarlens}[title={What you've built}]
One more word for telling, explaining and describing at B1.
\end{grammarlens}

\subsection*{Your turn}
\begin{itemize}
\raggedright
  \item \emph{Say it:} \emph{blank}
  \item \emph{Say it:} \emph{blank}, once more
  \item \emph{Say it:} a sentence of your own with \emph{blank}
  \item \emph{Recall:} say the Russian for a fringe, then the Russian for a certificate, then say \emph{blank} again
\end{itemize}

\begin{tcolorbox}[breakable,colback=teal!4,colframe=teal!35!black,title={Before you move on}]
What does \ru{бланк} mean? (\textquotedblleft{}A form (to fill in)\textquotedblright{}.) Say it once more.
\end{tcolorbox}
\section[\texorpdfstring{\ru{очередь} (óchered)}{очередь (óchered)}]{\ru{очередь} (óchered) — a queue; a turn}

\label{lesson:RU-C450-ochered}



\begin{quote}
Before the new one: say the Russian for a certificate, then the Russian for a form.
\end{quote}

\subsection*{You'll want to know: \ru{очередь}}
\textbf{\ru{очередь}} (\emph{óchered}) — \textquotedblleft{}a queue; a turn\textquotedblright{}.

\begin{quote}
\ru{На} \ru{почте} \ru{была} \ru{длинная} \ru{очередь}.
\end{quote}

\begin{quote}
\emph{Na póchte bylá dlínnaya óchered.}
\end{quote}

\begin{quote}
There was a long queue at the post office.
\end{quote}

Feminine. \emph{\ru{Стоять} \ru{в} \ru{очереди}} — to stand in a queue; \emph{\ru{моя} \ru{очередь}} — my turn.

\begin{cousinweb}
\emph{\ru{очередь}} is usually linked with \emph{\ru{череда}} 'a row, a turn in order'. EtymDB lists \emph{\ru{очередь}} without a source, so the line is hedged.
\end{cousinweb}

\begin{grammarlens}[title={What you've built}]
One more word for telling, explaining and describing at B1.
\end{grammarlens}

\subsection*{Your turn}
\begin{itemize}
\raggedright
  \item \emph{Say it:} \emph{óchered}
  \item \emph{Say it:} \emph{óchered}, once more
  \item \emph{Say it:} a sentence of your own with \emph{óchered}
  \item \emph{Recall:} say the Russian for a certificate, then the Russian for a form, then say \emph{óchered} again
\end{itemize}

\begin{tcolorbox}[breakable,colback=teal!4,colframe=teal!35!black,title={Before you move on}]
What does \ru{очередь} mean? (\textquotedblleft{}A queue; a turn\textquotedblright{}.) Say it once more.
\end{tcolorbox}
\section[\texorpdfstring{\ru{оформлять} (oformlyát)}{оформлять (oformlyát)}]{\ru{оформлять} (oformlyát) — to get (documents) done, to process; to design}

\label{lesson:RU-C450-oformlyat}



\begin{quote}
Before the new one: say the Russian for a form, then the Russian for a queue.
\end{quote}

\subsection*{You'll want to know: \ru{оформлять}}
\textbf{\ru{оформлять}} (\emph{oformlyát}) — \textquotedblleft{}to get (documents) done, to process; to design\textquotedblright{}.

\begin{quote}
\ru{Мы} \ru{три} \ru{дня} \ru{оформляли} \ru{визу}.
\end{quote}

\begin{quote}
\emph{My tri dnya oformlyáli vízu.}
\end{quote}

\begin{quote}
It took us three days to get the visa done.
\end{quote}

Imperfective; perfective \emph{\ru{оформить}} (\emph{\ru{оформлю}, \ru{оформишь}}).

\begin{cousinweb}
\emph{\ru{оформлять}} is built from \emph{\ru{о}-} and \emph{\ru{форма}} 'a form, a shape', which is usually traced to Latin \emph{forma} 'a shape' (Lewis \& Short; EtymDB lists \emph{\ru{форма}} without a source): to give a matter its proper form.
\end{cousinweb}

\begin{grammarlens}[title={What you've built}]
One more word for telling, explaining and describing at B1.
\end{grammarlens}

\subsection*{Your turn}
\begin{itemize}
\raggedright
  \item \emph{Say it:} \emph{oformlyát}
  \item \emph{Say it:} \emph{oformlyát}, once more
  \item \emph{Say it:} a sentence of your own with \emph{oformlyát}
  \item \emph{Recall:} say the Russian for a form, then the Russian for a queue, then say \emph{oformlyát} again
\end{itemize}

\begin{tcolorbox}[breakable,colback=teal!4,colframe=teal!35!black,title={Before you move on}]
What does \ru{оформлять} mean? (\textquotedblleft{}To get (documents) done, to process; to design\textquotedblright{}.) Say it once more.
\end{tcolorbox}
\section[\texorpdfstring{\ru{заверять} (zaveryát)}{заверять (zaveryát)}]{\ru{заверять} (zaveryát) — to certify; to assure}

\label{lesson:RU-C450-zaveryat}



\begin{quote}
Before the new one: say the Russian for a queue, then the Russian for to get.
\end{quote}

\subsection*{You'll want to know: \ru{заверять}}
\textbf{\ru{заверять}} (\emph{zaveryát}) — \textquotedblleft{}to certify; to assure\textquotedblright{}.

\begin{quote}
\ru{Нотариус} \ru{заверяет} \ru{копию} \ru{паспорта}.
\end{quote}

\begin{quote}
\emph{Notárius zaveryáyet kópiyu pásporta.}
\end{quote}

\begin{quote}
The notary certifies a copy of the passport.
\end{quote}

Imperfective; perfective \emph{\ru{заверить}}. \emph{\ru{Заверить} \ru{копию}} — to certify a copy.

\begin{cousinweb}
\emph{\ru{заверять}} is built from \emph{\ru{за}-} and \emph{-\ru{верять}}, from \emph{\ru{верить}} 'to believe', on \emph{\ru{вера}} 'faith'. EtymDB lists \emph{\ru{верить}} without a deeper source.
\end{cousinweb}

\begin{grammarlens}[title={What you've built}]
One more word for telling, explaining and describing at B1.
\end{grammarlens}

\subsection*{Your turn}
\begin{itemize}
\raggedright
  \item \emph{Say it:} \emph{zaveryát}
  \item \emph{Say it:} \emph{zaveryát}, once more
  \item \emph{Say it:} a sentence of your own with \emph{zaveryát}
  \item \emph{Recall:} say the Russian for a queue, then the Russian for to get, then say \emph{zaveryát} again
\end{itemize}

\begin{tcolorbox}[breakable,colback=teal!4,colframe=teal!35!black,title={Before you move on}]
What does \ru{заверять} mean? (\textquotedblleft{}To certify; to assure\textquotedblright{}.) Say it once more.
\end{tcolorbox}
\section[(dialogue)]{Words from chapters 443–450 — a review}

\label{lesson:RU-R450-review-443-450}



\begin{quote}
Say the Russian for to wander and to take a stroll.
\end{quote}

\subsection*{Your turn}
Say these aloud:

\begin{itemize}
\raggedright
  \item to wander, to take a stroll, a block, a suburb, a roadway
  \item a town hall, a cathedral, a bell tower, a fortress, a city hall
  \item a street lamp, a bench, a litter bin, a kiosk, a flower bed
  \item a small public garden, a shop sign, a shop window, a drainpipe, a stall
  \item a workshop, a shoemaker, a watchmaker, a locksmith, repairs
  \item a laundry, a dry cleaner's, to darn, to sew up, to sew on
  \item a hairdresser's, a haircut, to cut, to shave, a fringe
  \item a certificate, a form, a queue, to get, to certify
\end{itemize}

\begin{tcolorbox}[breakable,colback=teal!4,colframe=teal!35!black,title={Before you move on}]
To wander? (\textbf{\ru{бродить}}.) To take a stroll? (\textbf{\ru{прогуливаться}}.) A block? (\textbf{\ru{квартал}}.) A suburb? (\textbf{\ru{пригород}}.) A roadway? (\textbf{\ru{мостовая}}.) A town hall? (\textbf{\ru{ратуша}}.) A cathedral? (\textbf{\ru{собор}}.) A bell tower? (\textbf{\ru{колокольня}}.) A fortress? (\textbf{\ru{крепость}}.) A city hall? (\textbf{\ru{мэрия}}.) A street lamp? (\textbf{\ru{фонарь}}.) A bench? (\textbf{\ru{скамейка}}.) A litter bin? (\textbf{\ru{урна}}.) A kiosk? (\textbf{\ru{киоск}}.) A flower bed? (\textbf{\ru{клумба}}.) A small public garden? (\textbf{\ru{сквер}}.) A shop sign? (\textbf{\ru{вывеска}}.) A shop window? (\textbf{\ru{витрина}}.) A drainpipe? (\textbf{\ru{водосток}}.) A stall? (\textbf{\ru{ларёк}}.) A workshop? (\textbf{\ru{мастерская}}.) A shoemaker? (\textbf{\ru{сапожник}}.) A watchmaker? (\textbf{\ru{часовщик}}.) A locksmith? (\textbf{\ru{слесарь}}.) Repairs? (\textbf{\ru{ремонт}}.) A laundry? (\textbf{\ru{прачечная}}.) A dry cleaner's? (\textbf{\ru{химчистка}}.) To darn? (\textbf{\ru{штопать}}.) To sew up? (\textbf{\ru{зашивать}}.) To sew on? (\textbf{\ru{пришивать}}.) A hairdresser's? (\textbf{\ru{парикмахерская}}.) A haircut? (\textbf{\ru{стрижка}}.) To cut? (\textbf{\ru{стричь}}.) To shave? (\textbf{\ru{бриться}}.) A fringe? (\textbf{\ru{чёлка}}.) A certificate? (\textbf{\ru{справка}}.) A form? (\textbf{\ru{бланк}}.) A queue? (\textbf{\ru{очередь}}.) To get? (\textbf{\ru{оформлять}}.) To certify? (\textbf{\ru{заверять}}.)
\end{tcolorbox}
